\section{Conclusion and Outlook}
\label{sec:conclusion}

We presented \system, a persistent materials-research platform built on the \hermes agent runtime.  The system couples a controlled natural-language research environment to a structure-centred scientific service through a bounded MCP interface.  Five domain modules give the platform its scientific semantics: \matspec compiles open questions into constraint graphs; \matevidence aligns literature and federated database evidence around stable candidate identities; \mattransform realizes soft-chemistry proposals as typed structure operations or condition plans; \matdag constructs capability-aware scientific task graphs; and \matevolve carries candidate lineages and validation state into subsequent cycles.

The three case studies demonstrate the breadth of this formulation.  In the two-dimensional flat-band study, source-native database bands led directly to four explicit child structures and a structure-screened portfolio.  In the van der Waals ferromagnetic-topology study, the system organized six physically distinct routes under one five-observable verification matrix.  In the fractional-valence Lieb-lattice study, federated database discovery, literature-and-mechanism transfer, and structure generation converged into 21 hypotheses, 17 executed structures, six unique registry children, and a candidate-specific learned-Hamiltonian frontier.  The nearest current candidates---Li$_2$Ni$_5$OF$_{11}$ descendants and NbCl$_2$O/TaCl$_2$O oxyhalides---now have explicit next tasks for orbital character, charge distribution, lattice connectivity, crossing removal, and calibrated electronic confirmation.

These results show why a materials agent benefits from persistent scientific objects.  Evidence can remain incomplete without losing provenance; a proposed chemical change can become a reproducible child structure; a model result can update only the constraint it actually measures; and the next cycle can operate on the same candidate rather than reconstructing it from dialogue.  \hermes supplies interaction, policy, continuation, and workflow evolution, while the materials layer supplies the candidate, structure, evidence, and task semantics needed for scientific continuity.

The next release will extend this foundation in three directions.  First, structure generation will add ordered-supercell and partial-occupancy enumeration so that fractional substitution and mixed valence can be realized between end members.  Second, the electronic route will add orbital-resolved projections, contributor-sublattice graph validation, topological invariants, and calibrated uncertainty for learned Hamiltonians, followed by selected first-principles confirmation.  Third, the frozen layered benchmark will be executed on an expanded structure-family-independent case pool, enabling paired expert evaluation of the base model, \hermes harness, evidence layer, structure layer, and complete system.  These additions preserve the same candidate-centred contract and extend the set of scientific questions that can be resolved within it.

By placing evidence, structure manipulation, computation, and research evolution in one durable graph, \system provides a practical route from an open materials question to an explicit and progressively testable materials hypothesis.
